% Source: https://artofproblemsolving.com/community/c3h2530650_mathcounts_problem
% Stable identity: aops-topic-2530650-q01
Count the possible intersection points pair by pair:
\begin{itemize}
\item the two circles meet in at most $2$ points;
\item each of the $3\cdot2=6$ line--circle pairs meets in at most $2$ points, at most $12$ in all;
\item each of the $\dbinom32=3$ pairs of lines meets in at most $1$ point, at most $3$ in all.
\end{itemize}
So there are at most $2+12+3=17$ such points.

All $17$ can be distinct. For example, take the circles of radius $1.5$ centred at $(\pm0.8,0)$ and the lines
\[y=-0.9x-0.36,\qquad y=0.9x-0.55,\qquad y=-0.075x+0.63.\]
Each line passes within distance $1.5$ of both centres, so it cuts each circle twice; no two lines are parallel; and all the intersection points are different, as the sketch shows.

\begin{asy}
usepackage("amsmath");
size(6cm);
defaultpen(fontsize(10pt));
pair O1=(-0.8,0), O2=(0.8,0); real r=1.5;
path c1=circle(O1,r), c2=circle(O2,r);
real[] m={-0.9,0.9,-0.075}; real[] q={-0.36,-0.55,0.63};
path[] L;
for(int i=0;i<3;++i){ real x0=-2.75,x1=2.75; L.push((x0,m[i]*x0+q[i])--(x1,m[i]*x1+q[i])); }
draw(c1,linewidth(0.8)); draw(c2,linewidth(0.8));
for(int i=0;i<3;++i) draw(L[i],blue);
pair[] P;
P.append(intersectionpoints(c1,c2));
for(int i=0;i<3;++i){ P.append(intersectionpoints(L[i],c1)); P.append(intersectionpoints(L[i],c2)); }
for(int i=0;i<3;++i) for(int j=i+1;j<3;++j) P.append(intersectionpoints(L[i],L[j]));
for(int k=0;k<P.length;++k) dot(P[k],red);
label("$\ell_1$",(-2.75,-0.9*(-2.75)-0.36),W,blue);
label("$\ell_2$",(-2.75,0.9*(-2.75)-0.55),W,blue);
label("$\ell_3$",(2.75,-0.075*2.75+0.63),E,blue);
\end{asy}

The maximum is $\boxed{17}$.
